\documentclass[11pt, letterpaper]{article}
\usepackage{afp-style}
\afpdoctype{Party Constitution}

\begin{document}

\afptitlepage{Constitution of the\\Anti-Federalist Party}{The Foundational Governing Document}

\tableofcontents
\newpage

% ═══════════════════════════════════════════════════
\article{Name and Purpose}
% ═══════════════════════════════════════════════════

\subsection{Name}
This organization shall be known as \textbf{The Anti-Federalist Party} (hereinafter ``the Party'').

\subsection{Purpose}
The Party exists to restore governance to the level closest to the people. We hold that the concentration of political power in distant federal and state capitals has produced a government that no longer serves its citizens. Our purpose is to:

\begin{enumerate}
  \item Advance the principle that local self-governance is the only legitimate form of republican government.
  \item Elect candidates to local, state, and federal office who will decentralize power and return authority to communities.
  \item Educate the public about the Anti-Federalist intellectual tradition and its relevance to modern governance.
  \item Build a national network of autonomous local chapters that practice the decentralized governance they advocate.
  \item Defend individual liberty, data sovereignty, and community autonomy against all forms of centralized control --- governmental, corporate, or algorithmic.
\end{enumerate}

\afpcallout{``The government closest to the people serves the people best.'' This is not merely a slogan --- it is the organizing principle of this Party and every structure within it.}


% ═══════════════════════════════════════════════════
\article{Principles}
% ═══════════════════════════════════════════════════

The Anti-Federalist Party is founded upon the following immutable principles:

\subsection{Decentralization of Power}
Political authority derives from individuals and communities, not from federal or state governments. Power should be exercised at the most local level capable of effective governance.

\subsection{Individual Sovereignty}
Every person possesses inherent, inalienable rights that no government --- federal, state, or local --- may abridge. These include the rights enumerated in the Bill of Rights and the natural right to privacy, including digital privacy and data sovereignty.

\subsection{Community Self-Determination}
Counties, municipalities, and communities have the right to govern themselves according to their own values and priorities, provided they do not violate individual rights.

\subsection{Separation of Corporate and State Power}
The fusion of corporate and governmental power is as dangerous as the fusion of church and state. No corporation shall receive preferential treatment from any government the Party controls.

\subsection{Fiscal Responsibility}
No generation has the right to impose debt upon future generations without their consent. Government at all levels should operate within its means.

\subsection{Transparency and Accountability}
All government operations, expenditures, and decision-making processes must be transparent and accessible to the citizens they affect.

\subsection{Non-Interventionism}
The United States military exists to defend American soil. Foreign intervention, regime change, and permanent overseas military deployment are contrary to republican principles.


% ═══════════════════════════════════════════════════
\article{Membership}
% ═══════════════════════════════════════════════════

\subsection{Eligibility}
Membership in the Party is open to any person who:
\begin{enumerate}
  \item Is a citizen or lawful permanent resident of the United States, or a resident of a U.S.\ territory.
  \item Affirms the Principles set forth in Article~II of this Constitution.
  \item Completes a membership application and submits it to a recognized chapter or the national registry.
\end{enumerate}

\subsection{Dues}
\begin{enumerate}
  \item Membership dues, if any, shall be set by each chapter independently.
  \item No person shall be denied membership for inability to pay dues.
  \item The national organization shall not impose mandatory dues upon individual members.
\end{enumerate}

\subsection{Rights of Members}
Every member in good standing shall have the right to:
\begin{enumerate}
  \item Vote in chapter elections and on chapter resolutions.
  \item Run for any chapter or national office for which they are eligible.
  \item Attend and speak at any open meeting of their chapter or the national convention.
  \item Receive notice of all meetings, elections, and material decisions.
  \item Inspect the financial records of their chapter and the national organization.
\end{enumerate}

\subsection{Termination of Membership}
Membership may be terminated only by:
\begin{enumerate}
  \item Voluntary resignation submitted in writing.
  \item A two-thirds vote of the member's chapter, following notice and an opportunity to be heard, for conduct materially adverse to the Party's Principles.
\end{enumerate}


% ═══════════════════════════════════════════════════
\article{Chapter Structure}
% ═══════════════════════════════════════════════════

\subsection{Formation}
A chapter may be formed by any five (5) or more members within a defined geographic area (county, municipality, or equivalent) who:
\begin{enumerate}
  \item Adopt this Constitution and the Chapter Charter template.
  \item Elect initial officers as specified in Article~V.
  \item Submit a Charter Application to the national registry.
\end{enumerate}

\subsection{Autonomy}
Each chapter is an autonomous body. Chapters may:
\begin{enumerate}
  \item Set their own meeting schedules, dues, and operational procedures.
  \item Endorse candidates for local office independently.
  \item Pass resolutions on local, state, or national issues.
  \item Manage their own finances.
  \item Establish subcommittees as needed.
\end{enumerate}

\subsection{Obligations}
Each chapter shall:
\begin{enumerate}
  \item Operate in accordance with this Constitution and the Principles in Article~II.
  \item Hold a minimum of four (4) meetings per year.
  \item Maintain transparent financial records available to any member upon request.
  \item Report its membership count and officer roster to the national registry annually.
  \item Conduct elections for officers at least once every two (2) years.
\end{enumerate}

\subsection{Dissolution and Revocation}
\begin{enumerate}
  \item A chapter may dissolve itself by a two-thirds vote of its membership.
  \item A chapter's charter may be revoked by a three-fourths vote of the National Council for persistent, material violation of this Constitution, after notice and hearing.
  \item Upon dissolution, chapter assets shall be transferred to the state coordinating body or, in its absence, the national organization.
\end{enumerate}


% ═══════════════════════════════════════════════════
\article{Officers}
% ═══════════════════════════════════════════════════

\subsection{Chapter Officers}
Each chapter shall elect the following officers:
\begin{enumerate}
  \item \textbf{Chair} --- Presides over meetings, serves as the chapter's public representative, and coordinates with state and national bodies.
  \item \textbf{Vice Chair} --- Assumes the duties of the Chair in their absence and assists with chapter operations.
  \item \textbf{Secretary} --- Maintains meeting minutes, membership records, and official correspondence.
  \item \textbf{Treasurer} --- Manages chapter finances, maintains transparent records, and provides regular financial reports to the membership.
\end{enumerate}

\subsection{Terms of Office}
\begin{enumerate}
  \item All officers shall serve terms of two (2) years.
  \item No officer shall serve more than three (3) consecutive terms in the same position.
  \item Vacancies shall be filled by a special election within thirty (30) days.
\end{enumerate}

\subsection{Removal of Officers}
Any officer may be removed by a two-thirds vote of the chapter membership, following written notice of the specific charges and an opportunity for the officer to respond at a regular or special meeting.


% ═══════════════════════════════════════════════════
\article{National Structure}
% ═══════════════════════════════════════════════════

\afpcallout{The national organization exists to serve the chapters --- not to control them. The national structure is deliberately minimal. Power resides in the communities, not in a national headquarters.}

\subsection{National Council}
\begin{enumerate}
  \item The National Council shall consist of one (1) delegate from each chartered state coordinating body.
  \item The National Council shall meet at least twice per year, with at least one meeting conducted virtually to minimize travel burden.
  \item The National Council's authority is limited to: maintaining the national registry, coordinating national communications, organizing the national convention, and administering shared resources.
\end{enumerate}

\subsection{National Convention}
\begin{enumerate}
  \item A National Convention shall be held at least once every two (2) years.
  \item Each chapter shall be entitled to send delegates in proportion to its membership, with a minimum of one (1) delegate per chapter.
  \item The National Convention is the highest authority of the Party and may amend this Constitution.
\end{enumerate}

\subsection{State Coordinating Bodies}
\begin{enumerate}
  \item Chapters within a state may form a State Coordinating Body to coordinate statewide campaigns, share resources, and send a delegate to the National Council.
  \item State Coordinating Bodies have no authority over individual chapters.
  \item Formation requires a minimum of three (3) chapters within the state.
\end{enumerate}


% ═══════════════════════════════════════════════════
\article{Meetings and Decisions}
% ═══════════════════════════════════════════════════

\subsection{Regular Meetings}
Each chapter shall hold regular meetings at intervals determined by the chapter, but no fewer than four (4) per year.

\subsection{Special Meetings}
Special meetings may be called by the Chair or by petition of twenty percent (20\%) of chapter members.

\subsection{Quorum}
A quorum shall consist of one-third (1/3) of the chapter's membership in good standing, or five (5) members, whichever is greater.

\subsection{Voting}
\begin{enumerate}
  \item Each member present and in good standing shall have one (1) vote.
  \item Ordinary business shall be decided by simple majority vote.
  \item Constitutional amendments, charter revocations, officer removals, and dissolution require a two-thirds (2/3) supermajority.
  \item Proxy voting is not permitted. Virtual attendance with real-time participation shall be treated as present.
\end{enumerate}

\subsection{Transparency}
All meetings shall be open to the public unless the chapter votes to enter executive session for personnel or legal matters. Meeting minutes shall be made available to all members within fourteen (14) days.


% ═══════════════════════════════════════════════════
\article{Finances}
% ═══════════════════════════════════════════════════

\subsection{Revenue}
Chapter revenue may be derived from:
\begin{enumerate}
  \item Voluntary membership dues.
  \item Donations from individuals.
  \item Fundraising events.
  \item Grants from foundations aligned with the Party's Principles.
\end{enumerate}

\subsection{Prohibited Revenue}
The Party and its chapters shall not accept:
\begin{enumerate}
  \item Corporate donations or sponsorships.
  \item Donations from Political Action Committees (PACs) or Super PACs.
  \item Donations from foreign nationals or foreign governments.
  \item Any donation that creates or implies a quid pro quo obligation.
\end{enumerate}

\subsection{Financial Records}
\begin{enumerate}
  \item All financial records shall be maintained by the Treasurer and available for inspection by any member.
  \item The Treasurer shall present a financial report at each regular meeting.
  \item An annual financial summary shall be published to the full membership.
\end{enumerate}

\subsection{National Funding}
\begin{enumerate}
  \item The national organization may request voluntary contributions from chapters to fund shared resources (website, national convention, communications).
  \item No mandatory assessment shall be imposed upon chapters without a two-thirds vote of the National Convention.
\end{enumerate}


% ═══════════════════════════════════════════════════
\article{Candidate Endorsement}
% ═══════════════════════════════════════════════════

\subsection{Local Endorsements}
Each chapter may endorse candidates for local office by a two-thirds vote of members present at a regular or special meeting.

\subsection{State and Federal Endorsements}
\begin{enumerate}
  \item Endorsements for state office require approval by the State Coordinating Body, or by a majority of chapters within the state if no coordinating body exists.
  \item Endorsements for federal office require approval by the National Council or National Convention.
\end{enumerate}

\subsection{Endorsement Standards}
No candidate shall receive the Party's endorsement unless they publicly affirm:
\begin{enumerate}
  \item The Principles set forth in Article~II.
  \item A commitment to decentralize power if elected.
  \item Refusal to accept corporate PAC money.
\end{enumerate}


% ═══════════════════════════════════════════════════
\article{Amendments}
% ═══════════════════════════════════════════════════

\subsection{Proposal}
Amendments to this Constitution may be proposed by:
\begin{enumerate}
  \item A two-thirds vote of any chartered chapter.
  \item A majority vote of the National Council.
  \item A petition signed by ten percent (10\%) of the national membership.
\end{enumerate}

\subsection{Ratification}
Proposed amendments shall be ratified by a two-thirds vote at the National Convention, or by ratification votes in two-thirds of all chartered chapters within ninety (90) days of distribution.

\subsection{Immutable Principles}
The Principles enumerated in Article~II may not be amended or repealed. They are foundational and permanent. Any amendment that contradicts these Principles is void.


% ═══════════════════════════════════════════════════
\article{Dissolution}
% ═══════════════════════════════════════════════════

The Party may be dissolved only by a three-fourths vote of the National Convention, confirmed by ratification in three-fourths of all chartered chapters. Upon dissolution, all assets shall be distributed to nonpartisan civic education organizations selected by the final National Convention.


% ═══════════════════════════════════════════════════
\article{Ratification and Effective Date}
% ═══════════════════════════════════════════════════

This Constitution shall take effect upon adoption by the founding convention or, in the absence of a convention, upon ratification by the first five (5) chartered chapters.

\newpage

% ═══════════════════════════════════════════════════
% CERTIFICATION PAGE
% ═══════════════════════════════════════════════════

\thispagestyle{fancy}

\vspace*{2cm}

\afpcertification{Adopted by the Founding Convention}

\vspace{3cm}

\signatureline{Convention Chair}

\signatureline{Convention Secretary}

\signatureline{Presiding Officer}

\vspace{2cm}

\begin{center}
  {\small\color{afp-gray}\sffamily
    This Constitution was duly adopted by the founding members of the Anti-Federalist Party\\
    in accordance with the principles of decentralized, community-driven governance.\\[1em]
    All rights reserved. This document may be freely reproduced by any\\
    chartered chapter of the Anti-Federalist Party.
  }
\end{center}

\end{document}
